% ============================================================
%  I. Introduction
% ============================================================
\section{Introduction}

\IEEEPARstart{I}{n} the power generation, marine, and aerospace
industries, large rotating equipment---turbines, generators,
compressors---features long shafts whose installation accuracy
directly governs operational balance and service life.
%
Shaft alignment (or ``centreline finding'') is therefore a
critical step during installation and commissioning of such
machinery~\cite{piotrowski2006shaft}.

The prevailing industrial practice remains the
\emph{steel-wire method}, in which a taut piano wire is
stretched along the nominal shaft axis and the radial distance
from each bearing journal to the wire is measured with a
micrometer or dial indicator.
%
Despite its simplicity, the steel-wire method has well-known
drawbacks:
gravitational sag introduces systematic curvature
that must be compensated;
the wire is sensitive to air currents and mechanical vibration;
and the achievable measurement range is limited by wire tension
and sag correction accuracy.

Laser-based alignment offers a promising alternative.
%
By replacing the steel wire with a laser beam as the reference
axis, the problems of sag and vibration sensitivity are
eliminated.
%
However, the laser beam itself is invisible in free space;
its position can only be observed where it strikes a
receiving surface.
%
This paper proposes \textbf{DCPAM}
(Dual-Camera-based Point-to-Axis Measurement),
a model that employs two cameras---each observing the laser
spot on its own receiving screen---to reconstruct the
three-dimensional laser axis and compute the perpendicular
distance from an arbitrary target point to that axis.

The contributions of this work are as follows:
%
(1)~A complete mathematical formulation of the dual-camera
laser measurement pipeline, from pixel extraction to
point-to-axis distance computation.
%
(2)~Two parameter-tuning strategies:
DCPAM-CV (classical computer vision) and
DCPAM-CNN (convolutional neural network).
%
(3)~A closed-form error propagation analysis with
Monte Carlo validation.
%
(4)~A tilt-deviation rejection model that determines
the maximum allowable probe inclination for a given
target accuracy.

% TODO: 补充文献综述——现有激光对中方法、双目视觉测量、工业对中标准
